\section{Further research questions and concrete next experiments}
\label{sec:research}

The following program distinguishes short implementation projects,
structural mathematical questions, and obligations needed for a full
complexity theorem. A successful next step should produce either a proof
with explicit hypotheses, an independently replayable improvement, or an
informative counterexample.

\subsection{Immediate work on scalar transfers}

\begin{question}[Optimize the actual component cost]\label{q:order}
Can crossing orders and checkerboard colors be chosen to minimize merge
work, rather than only the global cut-edge or active-face maximum?
\end{question}

The current shading heuristic minimizes $(\max f_i,\sum f_i)$, while
Theorem~\ref{thm:componentwidth} depends on $g$ and actual work also depends
on alignment counts. A useful objective combines the sum of component orbit
counts, the products at merge steps, and estimated integer lengths.
Compute these structural features without running the full tensor
contraction, then compare greedy local changes on a fixed corpus.
The benchmark should include construction time and preserve orders on which
the existing heuristic wins. A worst-case approximation guarantee for this
objective would be more valuable than a favorable average alone.

\begin{question}[Switch representations at justified points]
Can a certified runtime policy choose among matching, one-tensor Potts,
and component Potts without paying their full combined costs?
\end{question}

The present measurements disprove a universal dominance claim. A first
policy can select an algorithm from cheap structural profiles and give each
attempt a fixed work budget. A stronger policy would convert a partially
computed representation when its cost grows. Such a conversion needs an
explicit linear map and an exact normalization account. Training a heuristic
on these 29 examples and reporting its accuracy on the same examples would
not answer the question; use held-out diagrams and relabeling controls.

\begin{question}[Factor beyond processed connected components]
When can an exact active tensor be split along articulation vertices,
small separators, or algebraic low-rank structure?
\end{question}

Graph-component independence is sufficient but not necessary for a compact
tensor. At a one-vertex separator, condition on the separator spin and
exploit the remaining color stabilizer. At larger separators, equality
patterns induce different stabilizer groups and orbit multiplicities.
The first milestone is a complete merge-and-forget calculus for separators
of size one and two, checked against labeled enumeration.
Algebraic rank factorization must preserve signed exact coefficients and
must charge for rank discovery. A large output rank supplies a natural
counterexample to an overly broad compression conjecture.

\begin{question}[Use contraction trees without assuming disk frontiers]
Can a tree-based exact Potts contraction improve path-order bottlenecks in
the repository while retaining the symmetry quotient?
\end{question}

Existing tensor and graph-width literature provides the starting point
\cite{meichanetzidis2019,maria2021}. The implementation question is how to
combine separator assignments, spin orbits and exact coefficient growth in
a reproducible cost model. A contraction tree should be stored as part of
the computation input and validated. Compare total contractions, maximum
tensor size, integer lengths and order-search cost. A tree decomposition
is useful only if its conversion and join operations are also bounded.

\subsection{Detection strength and exact specializations}

\begin{question}[Characterize the kernel of the six-color evaluation]
Which nontrivial knots, if any, have $V_K(2+\sqrt3)=1$?
\end{question}

This is a structural question about a specific evaluation. Start with
families having exact trace or skein recurrences: closed braids of fixed
index, twist constructions, satellites with controlled pattern, and
connected sums. Search for polynomial divisibility by $t^2-4t+1$ using
exact Laurent polynomials. A found collision is valuable even if another
filter already detects the knot. Failure to find one in a knot table gives
no completeness theorem.

\begin{question}[Design a small complementary menu]
Can a bounded or slowly growing set of algebraic points eliminate broad
classes of collisions while keeping total work small?
\end{question}

The $q=5$ failure and $q=6$ repair show why algebraic complementarity matters.
Measure additional detections per unit of exact work, with first-passing
and never-detected cases recorded. Distinguish changing $q$ from changing a
prime at fixed $q$. Proposition~\ref{prop:menu} gives an unconditional
growing menu for detecting whether the full Jones polynomial is $1$;
improving its number of points on a proved class is a concrete target.
A universal fixed menu would require a new theorem about the possible
Jones polynomials of knots.

\begin{question}[Arithmetic circuits or reconstruction]
Can modular evaluation and certified reconstruction outperform integer
pairs for very long, narrow diagrams?
\end{question}

The exact bit bound supplies an a priori reconstruction threshold.
One could evaluate over enough compatible primes, reconstruct each
coordinate, and verify with an additional exact identity. It is essential
to count prime selection, conversions, and total bit work. Pair arithmetic
has unusually cheap local multipliers, so a modular strategy may lose on
moderate inputs. Large-parameter tests should use the same diagram order
and report integer sizes, not merely crossing count.

\begin{question}[Shorter certificates than full replay]
Can a nontrivial exact Potts value be certified with a verifier materially
faster than recomputing the transfer?
\end{question}

The current evidence stores inputs and exact outputs, so replay has the
same complexity class as evaluation. A contraction transcript can expose
local sums and products, but may be as large as the computation.
Explore checkable arithmetic circuits, deterministic residual identities,
or interactive-verification models with their assumptions stated.
A soundness claim must account for arbitrary forged intermediate tables,
not just altered final values.

\subsection{Ring-valued compression for the complete fallback}

\begin{question}[Find useful ring-valued chain endomorphisms]
Can the scalar candidate search from the preceding continuation be extended
to produce small nontrivial same-matching ring-valued endomorphisms?
\end{question}

Theorem~\ref{thm:fitting} assumes the operator is supplied. A concrete
experiment is to add low dot-degree unknowns to the scalar intertwining
equations, solve them exactly, and test whether the lifted Fitting summands
reduce work on actual scanner blocks. Record the search dimension, solution
support, block sizes before and after splitting, and the cost of checking
all categorical cross terms. A search that routinely costs more than the
saved elimination is not a performance improvement.

\begin{question}[Exploit low annihilator degree]
When is the residue characteristic polynomial substantially larger than a
usable annihilator, and can a shorter one be found cheaply?
\end{question}

The universal certificate has degree $pr$, and the radical cutoff is sharp
in the full class. Particular matrices may have much smaller residue
minimal polynomials or additional relations. A residue minimal polynomial
alone need not lift with exponent $p$ by the proof of
Theorem~\ref{thm:residue}; that proof uses characteristic-polynomial
coefficients over the large ring. The first task is to identify a valid
lifting theorem or an explicit counterexample before replacing the
characteristic polynomial in code. Compare any successful refinement with
doubled geometric-series inversion.

\begin{question}[Keep Fitting projectors succinct]
Can $e(A)$ and its basis changes be manipulated without expanding their
coefficient supports or multiplicity spaces?
\end{question}

A degree-$O(r)$ polynomial is short, but evaluating it can create large
entries. Candidate representations include sparse monomial maps, shared
composition DAGs, and block recurrences. Each needs exact addition,
composition, zero testing, and application to differentials. The target
is a closure theorem bounding all these operations on a specified class.
Measure both the representation size and the cost of operations performed
on it; a small serialized expression can still hide expensive evaluation.

\begin{question}[Decision information without full homology]
Which exact summaries retain the rank-one question under all tangle
extensions while discarding irrelevant homology?
\end{question}

Previous saturated decisions and homological windows suggest useful
restricted summaries, but a complete invariant of future continuations can
be much larger than a scalar Euler characteristic. Formulate a congruence
on partial complexes that is respected by every allowed crossing extension.
Prove that the terminal decision descends to the quotient. Then either
bound the number of quotient states or exhibit a family requiring many.
This is a direct route toward the missing hypotheses of
Theorem~\ref{thm:conditional}.

\begin{question}[Learn from the explicit three-braid barrier]
Which features permit compressed three-braid algorithms, and which survive
at four or more strands?
\end{question}

The 2026 results provide both a warning about explicit output size and a
positive compressed model \cite{kelomaki2026}. A useful next project is to
encode a proven fixed-strand family with exact multiplicities and verify
the closed recurrences against the baseline scanner.
The research target is an extension theorem with a precise parameter,
not an empirical claim that all fixed-width complexes remain small.

\subsection{A fully costed geometric branch}

\begin{question}[Implement compressed normal-curve operations]
Can the published intersection and isotopy algorithms be integrated with
ProveIt's boundary-pattern representation?
\end{question}

The initial deliverable should handle explicit triangulated surfaces,
binary normal coordinates, geometric intersection numbers and the essential
curve precondition for isotopy \cite{lackenbycurves2026}.
Tests must include exponentially large weights described by short binary
inputs. Enumerating all normal arcs would destroy the intended complexity.
Every conversion between patterns, triangulations and normal coordinates
must be part of the time and size account.

\begin{question}[Control hierarchy restarts and encodings]
What invariant gives a useful simultaneous bound on hierarchy length,
pattern complexity, representation size and restart cost?
\end{question}

The expression $L(G+1)^L$ motivates a quantitative target but is not itself
a complete running-time theorem. Instrument an implementation of one
well-defined hierarchy step and prove that its compressed output meets a
stated bound. Then analyze how cutting and simplification interact with
that bound under restarts. A fast terminal checker is already a useful
component; the next milestone is a fully costed nonterminal transition,
followed by a global potential argument.

\begin{question}[Connect the algebraic and geometric branches]
Can an exact algebraic obstruction or partial homological summary guide
geometric choices without becoming an unproved pruning rule?
\end{question}

A safe interaction uses algebraic information to prioritize work while
retaining every branch required by the geometric correctness proof.
A stronger pruning rule needs its own implication theorem.
The first practical milestone is shared diagram and certificate formats
with explicit preprocessing equivalences. The longer-term milestone is
a combined algorithm whose worst-case bound can be proved independently
of its heuristic ordering choices.

\subsection{Evaluation discipline and recommended sequence}

Before a broad default change, extend the corpus beyond small named
fixtures and short random braids. Include pairs of diagrams of the same
knot with different embeddings and orders, nontrivial Alexander-trivial
knots, difficult unknot diagrams, large connected sums, and families chosen
to stress each representation. Record cases solved before the Jones stage
separately. Measure failed or budget-limited attempts as costs, not as
missing data.

The recommended sequence is: first improve the component-aware order
objective and confirm it on held-out inputs; next extract actual
same-matching blocks to test the residue certificates; then implement one
compressed geometric operation with a complete bit-complexity contract.
In parallel, maintain the exact-specialization collision search and the
independent oracles. This sequence produces useful verified increments
while keeping the eventual quasi-polynomial theorem tied to explicit,
testable obligations.
